
\subsection{$A_\infty$ 同态之间的同伦}
\label{subsec:4-6-2}
\liftIndex{I175}

本小节旨在证明：推论 4.6.3 所得滤过 $A_\infty$ 同态
$\mathfrak g$ 的同伦类不依赖于构造时的各种选择，
也不依赖于 $\psi$ 的辛同痕。换言之，我们要证明定理 4.1.2。
继续沿用上一小节的记号。为标明
$\mathfrak f$ 对 $\psi$、$\{J_\rho\}_\rho$
及 $\mathfrak s^{\operatorname{top}(\rho)}$ 的依赖，记为
$\mathfrak f^{\psi,\{J_\rho\}_\rho,
\mathfrak s^{\operatorname{top}(\rho)}}$。
上一小节构造了滤过 $A_\infty$ 同态
\begin{equation}
\begin{aligned}
\mathfrak f^{\psi,\{J_\rho\}_\rho,
\mathfrak s^{\operatorname{top}(\rho)}}\colon
&(C_1(L;\Lambda_{0,\mathrm{nov}}),
  \mathfrak m_k^{J,\mathfrak s})\\
&\longrightarrow
(C_3(L';\Lambda_{0,\mathrm{nov}}),
  \mathfrak m_k^{J',\mathfrak s''}).
\end{aligned}
\tag{4.6.21}\label{eq:4-6-21}
\end{equation}
把它与由自然包含 $C_2(L';\mathbb Q)\subset C_3(L';\mathbb Q)$
诱导的滤过 $A_\infty$ 同态
\[
(C_2(L';\Lambda_{0,\mathrm{nov}}),
 \mathfrak m_k^{J',\mathfrak s'})
\longrightarrow
(C_3(L';\Lambda_{0,\mathrm{nov}}),
 \mathfrak m_k^{J',\mathfrak s''})
\]
的同伦逆复合（参见 (4.6.2.2)），得到
\begin{equation}
\begin{aligned}
\mathfrak g^{\psi,\{J_\rho\}_\rho,
\mathfrak s^{\operatorname{top}(\rho)}}\colon
&(C_1(L;\Lambda_{0,\mathrm{nov}}),
  \mathfrak m_k^{J,\mathfrak s})\\
&\longrightarrow
(C_2(L';\Lambda_{0,\mathrm{nov}}),
  \mathfrak m_k^{J',\mathfrak s'}).
\end{aligned}
\tag{4.6.22}\label{eq:4-6-22}
\end{equation}
它是同伦等价。本小节以下把
$\mathfrak s''_{s=0}$、$\mathfrak s^{\operatorname{top}(\rho)}_{s=0}$
分别简记为 $\mathfrak s''$、$\mathfrak s^{\operatorname{top}(\rho)}$，
并写 $C_3^{s=0}(L';\Lambda_{0,\mathrm{nov}})$、
$\mathcal X_3^{s=0}(L')$。同样使用
$\mathfrak g_{s=0}^{\psi,\{J_\rho^{s=0}\}_\rho,
\mathfrak s^{\operatorname{top}(\rho)}_{s=0}}$
及对应的 $\mathfrak f_{s=0}$。

设 $\psi_s\colon M\to M'$（$s\in[0,1]$）是辛同痕，
满足 $\psi_0=\psi$、$\psi_1=\psi'$。
为清楚区分 $\psi_s$ 的源与靶，把靶流形记为 $M'$。
记整个族为 $\boldsymbol\psi=\{\psi_s\}_s$，并假设
\begin{equation}
\psi_s(L)=L'\qquad\text{对每个 }s\in[0,1].
\tag{4.6.23}\label{eq:4-6-23}
\end{equation}
把上一小节的构造应用于 $\psi_1=\psi'$（代替
$\psi_0=\psi$），得到：
\begin{enumerate}[label=\textup{(4.6.24.\arabic*)},leftmargin=5.9em]
 \item
 $\{J_\rho^{s=1}\}_\rho$，即 $M'$ 上从
 $J^{\psi_1}$ 到 $J'$ 的相容殆复结构同伦。
 \item
 光滑奇异链的可数集 $\mathcal X_3^{s=1}(L')$，
 它张成 $S(L')$ 的子复形 $C_3^{s=1}(L';\mathbb Q)$。
 \item
 多值截面系统
 $\mathfrak s''_{s=1}$ 与
 $\mathfrak s^{\operatorname{top}(\rho)}_{s=1}$；
 后者定义在
 \[
 \mathcal M_{k+1}^{\mathrm{main}}
 (M',L',\{J_\rho^{s=1}\}_\rho:
 \beta;\operatorname{top}(\rho);\vec P)
 \]
 上，
 前者定义在
 $\mathcal M_{k+1}^{\mathrm{main}}
 ((L',J'),\beta;\vec P')$ 上，
 其中 $P'_i\in\mathcal X_3^{s=1}(L')$。
\end{enumerate}
它们满足 (4.6.15)，因而诱导滤过 $A_\infty$ 同态
\[
\begin{aligned}
\mathfrak f_{s=1}^{\psi_1,\{J_\rho^{s=1}\}_\rho,
\mathfrak s^{\operatorname{top}(\rho)}_{s=1}}
&\colon(C_1(L;\Lambda_{0,\mathrm{nov}}),
\mathfrak m_k^{J,\mathfrak s})
\longrightarrow
(C_3^{s=1}(L';\Lambda_{0,\mathrm{nov}}),
\mathfrak m_k^{J',\mathfrak s''_{s=1}}),\\
\mathfrak g_{s=1}^{\psi_1,\{J_\rho^{s=1}\}_\rho,
\mathfrak s^{\operatorname{top}(\rho)}_{s=1}}
&\colon(C_1(L;\Lambda_{0,\mathrm{nov}}),
\mathfrak m_k^{J,\mathfrak s})
\longrightarrow
(C_2(L';\Lambda_{0,\mathrm{nov}}),
\mathfrak m_k^{J',\mathfrak s'}).
\end{aligned}
\]

现在给出本小节的主要结论。
\begin{theorem}{4.6.25}\label{thm:4-6-25}
$\mathfrak g_{s=0}^{\psi_0,\{J_\rho^{s=0}\}_\rho,
\mathfrak s^{\operatorname{top}(\rho)}_{s=0}}$
与
$\mathfrak g_{s=1}^{\psi_1,\{J_\rho^{s=1}\}_\rho,
\mathfrak s^{\operatorname{top}(\rho)}_{s=1}}$
同伦。
\end{theorem}

\begin{proof}
我们在相容殆复结构的两条同伦
$\{J_\rho^{s=0}\}_\rho$、
$\{J_\rho^{s=1}\}_\rho$ 之间再取一条同伦。
具体地，取双参数族 $\{J_{\rho,s}\}_{\rho,s}$，满足
\begin{enumerate}[label=\textup{(4.6.26.\arabic*)},leftmargin=5.9em]
 \item
 $\{J_{\rho,0}\}_\rho=\{J_\rho^{s=0}\}_\rho$。
 \item
 $\{J_{\rho,1}\}_\rho=\{J_\rho^{s=1}\}_\rho$。
 \item
 $J_{0,s}=(\psi_s)_*J$，对每个 $s\in[0,1]$。
 \item
 $J_{1,s}=J'$，对每个 $s\in[0,1]$。
\end{enumerate}
相容殆复结构空间的可缩性保证这样的
$\{J_{\rho,s}\}_{\rho,s}$ 存在；
以下始终用此记号表示该双参数族。

固定 $s$ 时，$\rho\mapsto J_{\rho,s}$ 是一族相容殆复结构，
记为 $\{J_{\rho,s}\}_\rho$。
由此得到模空间
$\mathcal M_{k+1}^{\mathrm{main}}
(M',L',\{J_{\rho,s}\}_\rho:
\beta;\operatorname{top}(\rho))$。定义
\begin{equation}
\begin{aligned}
&\mathcal M_{k+1}^{\mathrm{main}}
(M',L',\{J_{\rho,s}\}_{\rho,s}:
\beta;\operatorname{top}(\rho),\operatorname{twp}(s))\\
&\qquad=
\bigcup_{s\in[0,1]}\{s\}\times
\mathcal M_{k+1}^{\mathrm{main}}
(M',L',\{J_{\rho,s}\}_\rho:
\beta;\operatorname{top}(\rho)).
\end{aligned}
\tag{4.6.27}\label{eq:4-6-27}
\end{equation}
在该模空间上定义评价映射
\[
\begin{aligned}
ev_0&\colon\mathcal M_{k+1}^{\mathrm{main}}
(M',L',\{J_{\rho,s}\}_{\rho,s}:
\beta;\operatorname{top}(\rho),\operatorname{twp}(s))
\longrightarrow L',\\
ev_i&\colon\mathcal M_{k+1}^{\mathrm{main}}
(M',L',\{J_{\rho,s}\}_{\rho,s}:
\beta;\operatorname{top}(\rho),\operatorname{twp}(s))
\longrightarrow L,\quad i=1,\ldots,k,\\
ev_s&\colon\mathcal M_{k+1}^{\mathrm{main}}
(M',L',\{J_{\rho,s}\}_{\rho,s}:
\beta;\operatorname{top}(\rho),\operatorname{twp}(s))
\longrightarrow[0,1].
\end{aligned}
\]
具体地，对模空间元素
$((\Sigma,\vec z),(u_\alpha),(\rho_\alpha),s)$，
令
\begin{align}
ev_0((\Sigma,\vec z),(u_\alpha),(\rho_\alpha),s)
 &=u_{\alpha_0}(z_0),
 \tag{4.6.28.1}\label{eq:4-6-28-1}\\
ev_i((\Sigma,\vec z),(u_\alpha),(\rho_\alpha),s)
 &=\psi_s^{-1}(u_{\alpha_i}(z_i)),
 \tag{4.6.28.2}\label{eq:4-6-28-2}\\
ev_s((\Sigma,\vec z),(u_\alpha),(\rho_\alpha),s)
 &=s.
 \tag{4.6.28.3}\label{eq:4-6-28-3}
\end{align}
对 $P_i\in\mathcal X_1(L)$（$i=1,\ldots,k$），
用 $ev_i$ 定义
\begin{equation}
\begin{aligned}
&\mathcal M_{k+1}^{\mathrm{main}}
(M',L',\{J_{\rho,s}\}_{\rho,s}:
\beta;\operatorname{top}(\rho),\operatorname{twp}(s);\vec P)\\
&\quad=
\mathcal M_{k+1}^{\mathrm{main}}
(M',L',\{J_{\rho,s}\}_{\rho,s}:
\beta;\operatorname{top}(\rho),\operatorname{twp}(s))
{}_{(ev_1,\ldots,ev_k)}\!\times
\prod_{i=1}^kP_i .
\end{aligned}
\tag{4.6.29}\label{eq:4-6-29}
\end{equation}
定理 4.6.25 的证明思路是：利用这些模空间，
构造从 $(C_1(L;\Lambda_{0,\mathrm{nov}}),
\mathfrak m_k^{J,\mathfrak s})$
到 $[0,1]\times L'$ 上光滑奇异链模所定义的
滤过 $A_\infty$ 代数的同态。
为定义后一代数，还需要一些记号。

我们将对 $[0,1]\times L'$ 与 $L'$ 上奇异链复形的
适当子复形使用同态
$\operatorname{Eval}_{s=s_0}\colon
C([0,1]\times L')\to C(L')$。
它诱导的上同调同态
$H^*([0,1]\times L')\to H^*(L')$
与包含 $L'\cong\{s_0\}\times L'\subset[0,1]\times L'$
诱导的映射相同。
这里讨论的是上同调而非同调；虽然使用奇异链表示
上同调元素，但这是经由庞加莱对偶进行的。
具体地，上述同态的庞加莱对偶是
\[
\begin{aligned}
H_*([0,1]\times L';\{0,1\}\times L')
&\longrightarrow H_{*-1}(\{0,1\}\times L')\\
&\longrightarrow H_{*-1}(\{s_0\}\times L'),
\end{aligned}
\]
即连接同态再接投影到
$H_{*-1}(\{s_0\}\times L')$ 因子。
我们将在链层次定义
$C([0,1]\times L')$ 与 $C(L')$，
以便也在链层次定义这一连接同态。
以下用花体 $\mathcal P$ 表示
$[0,1]\times L'$ 上的奇异单形，
用正体 $P,P'$ 表示 $L,L'$ 上的奇异单形。

\begin{definition}{4.6.30}\label{def:4-6-30}
设 $\mathcal P=(|\mathcal P|,f)$ 是光滑奇异单形，
$f\colon|\mathcal P|\to[0,1]\times L'$。
若对 $s_0=0,1$ 都满足下列条件，则称
$\mathcal P$ 为\emph{适配奇异单形}：\liftIndex{I176}
\begin{enumerate}[label=\textup{(4.6.31.\arabic*)},leftmargin=5.9em]
 \item
 $f^{-1}(\{s_0\}\times L')\cap|\mathcal P|$
 或等于 $|\mathcal P|$，或为空，
 或仅由一个面（余维任意）组成。
 \item
 若
 $f^{-1}(\{0\}\times L')\cap|\mathcal P|\ne\varnothing$，
 则
 $|\mathcal P|\subset[0,1/3)\times L'$；
 若
 $f^{-1}(\{1\}\times L')\cap|\mathcal P|\ne\varnothing$，
 则
 $|\mathcal P|\subset(2/3,1]\times L'$。
 \item
 设 $\Delta_a^m$ 是 $|\mathcal P|$ 的一个面，
 $\Delta_b^{m'}$ 是包含 $\Delta_a^m$ 的另一个面，
 $x\in\Delta_a^m\subset f^{-1}(\{s_0\}\times L')$，
 而 $\Delta_b^{m'}$ 不包含在
 $f^{-1}(\{s_0\}\times L')$ 中。
 则对任意
 $N\in T_x\Delta_b^{m'}\setminus
 T_x(\partial\Delta_b^{m'})$，
 $f_*(N)$ 的 $[0,1]$ 分量不为零。
\end{enumerate}
\end{definition}

把 $[0,1]\times L'$ 上的 $m$ 维奇异单形看作次数
$d=1+\dim L'-m$ 的上链。
记 $S_+^d([0,1]\times L')$ 为次数 $d$
的适配上链所成的阿贝尔群；
容易看出它构成上链复形。
$S(\{0,1\}\times L')$ 是其子复形。
定义
\[
S([0,1]\times L')
=\frac{S_+([0,1]\times L')}
{S(\{0,1\}\times L')}.
\]
它由满足定义 4.6.30 的条件、
但不完全包含在 $\{0,1\}\times L'$
中的奇异单形自由生成。

如下定义
$\operatorname{Eval}_{s=s_0}\colon
S([0,1]\times L')\to S(L')$；
这里 $S(L')$ 是 $L'$ 的光滑奇异链复形。
若 $f^{-1}(\{s_0\}\times L')\cap|\mathcal P|$
在 $|\mathcal P|$ 中不是余维一，则令
\begin{equation}
\operatorname{Eval}_{s=s_0}(\mathcal P)=0.
\tag{4.6.32.1}\label{eq:4-6-32-1}
\end{equation}
若它是余维一，则令
\begin{equation}
\operatorname{Eval}_{s=s_0}(\mathcal P)
=\pm(f^{-1}(\{s_0\}\times L')\cap|\mathcal P|,f).
\tag{4.6.32.2}\label{eq:4-6-32-2}
\end{equation}
符号 $\pm$ 见第 7.2.10 小节。
我们将在第 7.2 节的引理 7.2.136 证明，
$\operatorname{Eval}_{s=s_0}$ 是上链映射，
且是上链同伦等价。

对 $\mathcal P\in S([0,1]\times L')$
还考虑以下条件：
\begin{statement}{条件}{4.6.33}\label{cond:4-6-33}
$\operatorname{Eval}_{s=s_0}(\mathcal P)
\in C_3^{s=s_0}(L';\mathbb Q)$。
\end{statement}

按下式定义
$\operatorname{Incl}_{1,\beta_0}\colon
S(L')\to S([0,1]\times L')$：
\[
\operatorname{Incl}_{1,\beta_0}(|P|,f)
=([0,1]\times|P|,\operatorname{id}\times f).
\]
这里对 $[0,1]\times|P|$ 取适当的
棱柱型单纯剖分，把右端看作奇异链；
详见第 7.2.10 小节。

我们将在一个由可数多个适配光滑奇异单形张成的
上链复形上构造滤过 $A_\infty$ 代数。
为此使用伪全纯圆盘模空间的另一个参数化版本。
对满足 (4.6.26) 的族
$\{J_{\rho,s}\}_{\rho,s}$，令 $\rho=1$，
得到相容殆复结构族
$\{J_{1,s}\}_s$（$s\in[0,1]$）。
由 (4.6.26.4)，这是常值族：
对每个 $s$，$J_{1,s}=J'$。
从它得到模空间
$\mathcal M_{k+1}^{\mathrm{main}}((L',J_{1,s}),\beta)$。
定义\liftIndex{I177}
\begin{equation}
\begin{aligned}
&\mathcal M_{k+1}^{\mathrm{main}}
(M',L',\{J_{1,s}\}_s:\beta;\operatorname{twp}(s))\\
&\quad=
\bigcup_{s\in[0,1]}\{s\}\times
\mathcal M_{k+1}^{\mathrm{main}}((L',J_{1,s}),\beta)
=[0,1]\times
\mathcal M_{k+1}^{\mathrm{main}}((L',J'),\beta).
\end{aligned}
\tag{4.6.34}\label{eq:4-6-34}
\end{equation}
这里以模空间
$\mathcal M_{k+1}^{\mathrm{main}}((L',J_{1,s}),\beta)$
的\emph{逐时积}定义
$\mathcal M_{k+1}^{\mathrm{main}}
(M',L',\{J_{1,s}\}_s:\beta;\operatorname{twp}(s))$。
这不同于使用族 $s\mapsto J_{1,s}$
构造的
$\mathcal M_{k+1}^{\mathrm{main}}
(M',L',\{J_{1,s}\}_s:
\beta;\operatorname{top}(s))$：
后者是半稳定带边黎曼曲面上的伪全纯映射模空间，
殆复结构可随分支而变；
前者的元素则对应某个固定的 $s$，
在各分支上都使用同一个 $J_{1,s}$。

定义评价映射
\[
\begin{aligned}
ev_i\colon\mathcal M_{k+1}^{\mathrm{main}}
(M',L',\{J_{1,s}\}_s:\beta;\operatorname{twp}(s))
&\longrightarrow L',\quad i=0,\ldots,k,\\
ev_s\colon\mathcal M_{k+1}^{\mathrm{main}}
(M',L',\{J_{1,s}\}_s:\beta;\operatorname{twp}(s))
&\longrightarrow[0,1]
\end{aligned}
\]
如下：
\begin{align}
ev_i((\Sigma,\vec z),(u_\alpha),s)
 &=u_{\alpha_i}(z_i),
 \tag{4.6.35.1}\label{eq:4-6-35-1}\\
ev_s((\Sigma,\vec z),(u_\alpha),s)
 &=s.
 \tag{4.6.35.2}\label{eq:4-6-35-2}
\end{align}
注意这里每个 $ev_i$ 都映到 $L'$。
记 $ev_i^+=(ev_s,ev_i)$。
对 $\mathcal P_i\in S([0,1]\times L')$
（$i=1,\ldots,k$），定义
\[
\begin{aligned}
&\mathcal M_{k+1}^{\mathrm{main}}
(M',L',\{J_{1,s}\}_s:
\beta;\operatorname{twp}(s);\vec{\mathcal P})\\
&\quad=
\mathcal M_{k+1}^{\mathrm{main}}
(M',L',\{J_{1,s}\}_s:
\beta;\operatorname{twp}(s))
{}_{(ev_1^+,\ldots,ev_k^+)}\!\times
\prod_{i=1}^k\mathcal P_i .
\end{aligned}
\]
这里的纤维积取在 $([0,1]\times L')^k$ 上。
此空间有仓西结构。
若选取与零横截的多值截面
$\mathfrak s^{\operatorname{twp}(s)}$，
则得到扰动后的模空间
\[
\mathcal M_{k+1}^{\mathrm{main}}
(M',L',\{J_{1,s}\}_s:
\beta;\operatorname{twp}(s);\vec{\mathcal P})
^{\mathfrak s}
=((\mathfrak s^{\operatorname{twp}(s)})^{-1}(0),ev_0^+).
\]
可以对其零集作单纯剖分，
使右端成为 $[0,1]\times L'$ 中的奇异链，
且每个单形都适配；见第 7.2.10 小节。
并有
\begin{equation}
\begin{aligned}
&\deg\mathcal M_{k+1}^{\mathrm{main}}
(M',L',\{J_{1,s}\}_s:
\beta;\operatorname{twp}(s);\vec{\mathcal P})
^{\mathfrak s}\\
&\qquad=2-k+\sum_i\deg\mathcal P_i-\mu(\beta).
\end{aligned}
\tag{4.6.36.1}\label{eq:4-6-36-1}
\end{equation}
这里 $\deg$ 是链的余维，即上同调次数。
上述模空间的边界也可表为同类模空间的纤维积。

另一方面，在
$\mathcal M_{k+1}^{\mathrm{main}}
(M',L',\{J_{\rho,s}\}_{\rho,s}:
\beta;\operatorname{top}(\rho),\operatorname{twp}(s);
\vec P)$
上取与零横截的多值截面
$\mathfrak s^{\operatorname{top}(\rho),\operatorname{twp}(s)}$，
同样得到扰动后的模空间
\[
\begin{aligned}
&\mathcal M_{k+1}^{\mathrm{main}}
(M',L',\{J_{\rho,s}\}_{\rho,s}:
\beta;\operatorname{top}(\rho),\operatorname{twp}(s);
\vec P)^{\mathfrak s}\\
&\quad=
((\mathfrak s^{\operatorname{top}(\rho),
\operatorname{twp}(s)})^{-1}(0),ev_0^+).
\end{aligned}
\]
可对该双参数截面的零集作单纯剖分，使该对成为
$[0,1]\times L'$ 中的奇异链，
且每个单形满足定义 4.6.30 的条件；
见第 7.2.10 小节。\translatornote{原书此处仍写
$(\mathfrak s^{\operatorname{twp}(s)})^{-1}(0)$，
但本段所用的截面是
$\mathfrak s^{\operatorname{top}(\rho),
\operatorname{twp}(s)}$；
此处按正在讨论的双参数模空间定型，
详见疑误 A-0242。}
其次数为
\begin{equation}
\begin{aligned}
&\deg\mathcal M_{k+1}^{\mathrm{main}}
(M',L',\{J_{\rho,s}\}_{\rho,s}:
\beta;\operatorname{top}(\rho),
\operatorname{twp}(s);\vec P)^{\mathfrak s}\\
&\qquad=3-k+\sum_i\deg P_i-\mu(\beta).
\end{aligned}
\tag{4.6.36.2}\label{eq:4-6-36-2}
\end{equation}

以下命题是我们在 $S([0,1]\times L';\partial)$
（即上述相对商复形）
的可数生成子复形上构造滤过 $A_\infty$ 结构
所需的主要结果。
其证明将在第 7.2.10 小节给出。
严格地说，那里先证明命题 4.6.37 的
$A_{n,K}$ 版本（即命题 7.2.160、7.2.162），
再于第 7.2.11--7.2.13 小节推出 $A_\infty$ 版本。

\begin{proposition}{4.6.37}\label{prop:4-6-37}
存在 $[0,1]\times L'$ 上奇异单形的可数集
$\mathcal X([0,1]\times L')$，
它张成 $S([0,1]\times L')$ 的子复形
$C([0,1]\times L')$；
还存在与零横截的多值截面系统
$\mathfrak s^{\operatorname{twp}(s)}$、
$\mathfrak s^{\operatorname{top}(\rho),\operatorname{twp}(s)}$，
满足以下性质：
\begin{enumerate}[label=\textup{(4.6.38.\arabic*)},leftmargin=5.9em]
 \item
 $\mathcal X([0,1]\times L')$ 的每个元素均适配，
 且满足条件 4.6.33；
 它们均不完全包含在 $\{0,1\}\times L'$ 中。
 \item
 若 $\mathcal P_i\in\mathcal X([0,1]\times L')$，
 则
 \[
 \mathcal M_{k+1}^{\mathrm{main}}
 (M',L',\{J_{1,s}\}_s:
 \beta;\operatorname{twp}(s);\vec{\mathcal P})
 ^{\mathfrak s}
 \in C([0,1]\times L').
 \]
 \item
 $\mathfrak s^{\operatorname{twp}(s)}$
 与
 $\mathfrak s^{\operatorname{top}(\rho),\operatorname{twp}(s)}$
 在各自模空间的边界上相容。
 \item
 若 $P\in\mathcal X_2(L')$，
 则 $\operatorname{Incl}_{1,\beta_0}(P)
 \in C([0,1]\times L')$。
 \item
 包含映射
 $C([0,1]\times L')\to S([0,1]\times L')$
 在上同调群上诱导同构。
 \item
 对 $\mathcal P_i\in C([0,1]\times L')$，
 记 $\vec{\mathcal P}=(\mathcal P_1,\ldots,\mathcal P_k)$，
 并令
 $\mathcal P_i|_{s=0}
 =\operatorname{Eval}_{s=0}(\mathcal P_i)$、
 $\vec{\mathcal P}|_{s=0}
 =(\mathcal P_1|_{s=0},\ldots,\mathcal P_k|_{s=0})$。
 则
 \[
 \begin{aligned}
 &\operatorname{Eval}_{s=0}\!
 \bigl(\mathcal M_{k+1}^{\mathrm{main}}
 (M',L',\{J_{1,s}\}_s:
 \beta;\operatorname{twp}(s);
 \vec{\mathcal P})^{\mathfrak s}\bigr)\\
 &\qquad=
 \mathcal M_{k+1}^{\mathrm{main}}
 ((M',L'),J':\beta;
 \vec{\mathcal P}|_{s=0})^{\mathfrak s''}.
 \end{aligned}
 \]
 同样令
 $\mathcal P_i|_{s=1}
 =\operatorname{Eval}_{s=1}(\mathcal P_i)$、
 $\vec{\mathcal P}|_{s=1}
 =(\mathcal P_1|_{s=1},\ldots,\mathcal P_k|_{s=1})$，
 则
 \[
 \begin{aligned}
 &\operatorname{Eval}_{s=1}\!
 \bigl(\mathcal M_{k+1}^{\mathrm{main}}
 (M',L',\{J_{1,s}\}_s:
 \beta;\operatorname{twp}(s);
 \vec{\mathcal P})^{\mathfrak s}\bigr)\\
 &\qquad=
 \mathcal M_{k+1}^{\mathrm{main}}
 ((M',L'),J':\beta;
 \vec{\mathcal P}|_{s=1})^{\mathfrak s''_{s=1}}.
 \end{aligned}
 \]
 这里 (4.6.26.4) 保证 $J_{1,0}=J_{1,1}=J'$。
 \item
 对 $P_i\in C(L)$，有
 \[
 \mathcal M_{k+1}^{\mathrm{main}}
 (M',L',\{J_{\rho,s}\}_{\rho,s}:
 \beta;\operatorname{top}(\rho),
 \operatorname{twp}(s);\vec P)^{\mathfrak s}
 \in C([0,1]\times L').
 \]
 \item
 可以定义
 \[
 \operatorname{Incl}_{k,\beta}\colon
 B_k(C(L')[1])\longrightarrow
 C([0,1]\times L')
 \]
 使 $\operatorname{Incl}_{1,\beta_0}$
 与 $\operatorname{Incl}_{k,\beta}$ 一起
 给出滤过 $A_\infty$ 同态
 $\operatorname{Incl}\colon
 C(L')[1]\to C([0,1]\times L')$，
 且满足
 $\operatorname{Eval}_{s=s_0}\circ
 \operatorname{Incl}=\operatorname{id}$。
\end{enumerate}
\end{proposition}

我们将用命题 4.6.37 构造图式 4.6.1。
图中以 $\mathfrak f_0$、$\mathfrak f_1$
及 $\mathfrak f_{\mathrm{tw}}$
分别缩写两端的同态及逐时积族给出的同态；
所有靶复形都在 $L'$ 上。\translatornote{原书图式
4.6.1 将两端的 $C_3^{s=0,1}$、$C_2$
和中间的区间复形都写在无撇的 $L$ 上；
但本段定义、(4.6.21)--(4.6.22)
及端点评价均以 $L'$ 为靶。
此处按已定义的 $C([0,1]\times L')$ 定型，
详见疑误 A-0241。}
\begin{figure}[htbp]
 \centering
 \setlength{\unitlength}{0.8mm}
 \begin{picture}(135,88)
  \put(15,43){\makebox(0,0){$C_1(L)$}}
  \put(83,78){\makebox(0,0){$C_3^{s=1}(L')$}}
  \put(83,43){\makebox(0,0){$C([0,1]\times L')$}}
  \put(83,8){\makebox(0,0){$C_3^{s=0}(L')$}}
  \put(125,78){\makebox(0,0){$\supset C_2(L')$}}
  \put(125,8){\makebox(0,0){$\supset C_2(L')$}}
  \put(27,44){\vector(1,0){42}}
  \put(44,47){\makebox(0,0)[b]{$\scriptstyle\mathfrak f_{\mathrm{tw}}$}}
  \put(25,48){\vector(2,1){42}}
  \put(40,64){\makebox(0,0)[b]{$\scriptstyle\mathfrak f_1$}}
  \put(25,38){\vector(2,-1){42}}
  \put(40,22){\makebox(0,0)[t]{$\scriptstyle\mathfrak f_0$}}
  \put(83,53){\vector(0,1){17}}
  \put(88,61){\makebox(0,0)[l]{$\scriptstyle\operatorname{Eval}_{s=1}$}}
  \put(83,34){\vector(0,-1){17}}
  \put(88,25){\makebox(0,0)[l]{$\scriptstyle\operatorname{Eval}_{s=0}$}}
  \put(122,35){\makebox(0,0){$\Vert$}}
  \put(122,51){\makebox(0,0){$\Vert$}}
 \end{picture}
 \par\smallskip
 \renewcommand{\thefigure}{4.6.1}
 \refstepcounter{figure}
 \textbf{图式~\thefigure}\label{diag:4-6-1}
\end{figure}

命题 4.6.37 将在第 7.2.10 小节证明。
注意该命题中的 $\operatorname{Incl}$
不满足 (4.2.3.1)；见注 7.2.172。

由此命题，在
$C([0,1]\times L';\Lambda_{0,\mathrm{nov}})$
上定义滤过 $A_\infty$ 代数结构。
这里的 $C([0,1]\times L';\Lambda_{0,\mathrm{nov}})$
是 $C([0,1]\times L')\otimes
\Lambda_{0,\mathrm{nov}}$ 的完备化。
设 $\mathcal P_i\in C([0,1]\times L')$，
当 $(k,\beta)\ne(1,\beta_0)=(1,0)$ 时，定义
\[
\mathfrak m_{k,\beta}
(\mathcal P_1,\ldots,\mathcal P_k)
=\mathcal M_{k+1}^{\mathrm{main}}
(M',L',\{J_{1,s}\}_s:
\beta;\operatorname{twp}(s);
\vec{\mathcal P})^{\mathfrak s}
\in C([0,1]\times L'),
\]
并令
$\mathfrak m_{1,\beta_0}(\mathcal P)
=(-1)^{n+1}\partial\mathcal P$，
其中 $n=\dim L$；
这里 $\mathcal P$ 是 $(n+1)$ 维空间
$[0,1]\times L'$ 上的奇异单形。
再令
\[
\mathfrak m_k
=\sum_\beta\mathfrak m_{k,\beta}
\otimes T^{\omega(\beta)}e^{\mu(\beta)/2}.
\]

\begin{lemma}{4.6.39}\label{lem:4-6-39}
$(C([0,1]\times L';\Lambda_{0,\mathrm{nov}}),
\mathfrak m)$ 是滤过 $A_\infty$ 代数。
\end{lemma}
证明与定理 3.5.11 的证明类似，
使用命题 4.6.37 即可，故从略。
接下来由 (4.6.38.4)、(4.6.38.6) 与 (4.6.38.7)
构造几个 $A_\infty$ 同态。首先定义
\[
\begin{aligned}
\operatorname{Eval}_{s=0}\colon
&(C([0,1]\times L';\Lambda_{0,\mathrm{nov}}),
  \mathfrak m)
\longrightarrow
(C_3^{s=0}(L';\Lambda_{0,\mathrm{nov}}),
  \mathfrak m_k^{J',\mathfrak s''_{s=0}}),\\
\operatorname{Eval}_{s=1}\colon
&(C([0,1]\times L';\Lambda_{0,\mathrm{nov}}),
  \mathfrak m)
\longrightarrow
(C_3^{s=1}(L';\Lambda_{0,\mathrm{nov}}),
  \mathfrak m_k^{J',\mathfrak s''_{s=1}}).
\end{aligned}
\]
其一次分量由 (4.6.32) 给出，
而 $k\ne1$ 时令
$(\operatorname{Eval}_{s=s_0})_k=0$。
由 (4.6.38.6)，
$\operatorname{Eval}_{s=s_0}$
是滤过 $A_\infty$ 同态。

再由 (4.6.38.8) 定义滤过 $A_\infty$ 同态
\[
\begin{aligned}
\operatorname{Incl}\colon
&(C_2(L';\Lambda_{0,\mathrm{nov}}),
  \mathfrak m_k^{J',\mathfrak s'})
\longrightarrow
(C([0,1]\times L';\Lambda_{0,\mathrm{nov}}),
  \mathfrak m).
\end{aligned}
\]
还可选择 $\mathcal X([0,1]\times L')$，
使前述
$(\operatorname{Eval}_{s=0})_1
\oplus(\operatorname{Eval}_{s=1})_1$
是满射；这一事实将在第 7.2.13 小节使用。
于是
$(C([0,1]\times L';\Lambda_{0,\mathrm{nov}}),
\mathfrak m)$
在某种广义意义下是
$[0,1]\times
(C(L';\Lambda_{0,\mathrm{nov}}),\mathfrak m_k)$
的模型；见定义 7.2.174。

此外，定义
\[
\begin{aligned}
\mathfrak f_{\mathrm{tw}}
:=
\mathfrak f^{\boldsymbol\psi,
\{J_{\rho,s}\}_{\rho,s},
\mathfrak s^{\operatorname{top}(\rho),
\operatorname{twp}(s)}}
\colon
&(C_1(L;\Lambda_{0,\mathrm{nov}}),
  \mathfrak m_k^{J,\mathfrak s})\\
&\longrightarrow
(C([0,1]\times L')
\widehat{\otimes}\Lambda_{0,\mathrm{nov}},
  \mathfrak m),
\end{aligned}
\]
其中
\[
\begin{aligned}
(\mathfrak f_{\mathrm{tw}})_k
&=\sum_\beta
  (\mathfrak f_{\mathrm{tw}})_{k,\beta}\,
  T^{\omega(\beta)}e^{\mu_L(\beta)/2},\\
(\mathfrak f_{\mathrm{tw}})_{k,\beta}
(P_1,\ldots,P_k)
&=
\mathcal M_{k+1}^{\mathrm{main}}
(M',L',\{J_{\rho,s}\}_{\rho,s}:
\beta;\operatorname{top}(\rho),
\operatorname{twp}(s);\vec P)^{\mathfrak s}.
\end{aligned}
\]
见图式 4.6.1。这里右端利用 (4.6.29)
及多值截面
$\mathfrak s^{\operatorname{top}(\rho),
\operatorname{twp}(s)}$ 定义；
由 (4.6.38.7)，它属于
$C([0,1]\times L')$。
与定理 4.6.1 的证明类似，
可证 $\mathfrak f_{\mathrm{tw}}$
是滤过 $A_\infty$ 同态。
下一引理由定义立即得到。

\begin{lemma}{4.6.40}\label{lem:4-6-40}
有
\begin{align}
\operatorname{Eval}_{s=0}\circ
\mathfrak f_{\mathrm{tw}}
&=
\mathfrak f_{s=0}^{\psi_0,
\{J_\rho^{s=0}\}_\rho,
\mathfrak s^{\operatorname{top}(\rho)}_{s=0}},
\tag{4.6.41.1}\label{eq:4-6-41-1}\\
\operatorname{Eval}_{s=1}\circ
\mathfrak f_{\mathrm{tw}}
&=
\mathfrak f_{s=1}^{\psi_1,
\{J_\rho^{s=1}\}_\rho,
\mathfrak s^{\operatorname{top}(\rho)}_{s=1}}.
\tag{4.6.41.2}\label{eq:4-6-41-2}
\end{align}
\end{lemma}

令
\[
\begin{aligned}
I_0\colon
&(C_2(L';\Lambda_{0,\mathrm{nov}}),
  \mathfrak m_k^{J',\mathfrak s'})
\longrightarrow
(C_3^{s=0}(L';\Lambda_{0,\mathrm{nov}}),
  \mathfrak m_k^{J',\mathfrak s''_{s=0}}),\\
I_1\colon
&(C_2(L';\Lambda_{0,\mathrm{nov}}),
  \mathfrak m_k^{J',\mathfrak s'})
\longrightarrow
(C_3^{s=1}(L';\Lambda_{0,\mathrm{nov}}),
  \mathfrak m_k^{J',\mathfrak s''_{s=1}})
\end{aligned}
\]
为包含同态，
$K_0,K_1$ 分别为其同伦逆。
由 (4.6.38.8) 还立即得到：
\begin{lemma}{4.6.42}\label{lem:4-6-42}
$\operatorname{Eval}_{s=0}\circ
\operatorname{Incl}=I_0$，
$\operatorname{Eval}_{s=1}\circ
\operatorname{Incl}=I_1$。
\end{lemma}
最后证明：
\begin{lemma}{4.6.43}\label{lem:4-6-43}
$K_0\circ\operatorname{Eval}_{s=0}$
与 $K_1\circ\operatorname{Eval}_{s=1}$ 同伦。
\end{lemma}
\begin{proof}
有
\[
\begin{aligned}
K_0\circ\operatorname{Eval}_{s=0}
\circ\operatorname{Incl}
&=K_0\circ I_0
\sim\operatorname{id}\\
&\sim K_1\circ I_1
=K_1\circ\operatorname{Eval}_{s=1}
\circ\operatorname{Incl}.
\end{aligned}
\]
由 (4.6.38) 与定理 4.2.45，
$\operatorname{Incl}$ 是同伦等价，
故结论成立。
\end{proof}

这些映射及端点关系可汇成图式 4.6.2。
为保持图中标记清晰，
用 $\mathfrak f_0,\mathfrak f_1$
缩写引理 4.6.40 右端的同态。
\begin{figure}[htbp]
 \centering
 \setlength{\unitlength}{0.8mm}
 \begin{picture}(140,80)
  \put(12,70){\makebox(0,0){$C_1$}}
  \put(12,40){\makebox(0,0){$C_1$}}
  \put(12,10){\makebox(0,0){$C_1$}}
  \put(73,70){\makebox(0,0){$C_3^{s=0}$}}
  \put(73,40){\makebox(0,0){$C([0,1]\times L')$}}
  \put(73,10){\makebox(0,0){$C_3^{s=1}$}}
  \put(128,70){\makebox(0,0){$C_2$}}
  \put(128,40){\makebox(0,0){$C_2$}}
  \put(128,10){\makebox(0,0){$C_2$}}
  \put(20,70){\vector(1,0){39}}
  \put(40,73){\makebox(0,0)[b]{$\scriptstyle\mathfrak f_0$}}
  \put(20,40){\vector(1,0){31}}
  \put(37,43){\makebox(0,0)[b]{$\scriptstyle\mathfrak f_{\mathrm{tw}}$}}
  \put(20,10){\vector(1,0){39}}
  \put(40,13){\makebox(0,0)[b]{$\scriptstyle\mathfrak f_1$}}
  \put(118,70){\vector(-1,0){30}}
  \put(103,73){\makebox(0,0)[b]{$\scriptstyle I_0$}}
  \put(118,40){\vector(-1,0){24}}
  \put(105,43){\makebox(0,0)[b]{$\scriptstyle\operatorname{Incl}$}}
  \put(118,10){\vector(-1,0){30}}
  \put(103,13){\makebox(0,0)[b]{$\scriptstyle I_1$}}
  \put(73,50){\vector(0,1){12}}
  \put(78,56){\makebox(0,0)[l]{$\scriptstyle\operatorname{Eval}_{s=0}$}}
  \put(73,30){\vector(0,-1){12}}
  \put(78,24){\makebox(0,0)[l]{$\scriptstyle\operatorname{Eval}_{s=1}$}}
  \put(12,55){\makebox(0,0){$\Vert$}}
  \put(12,25){\makebox(0,0){$\Vert$}}
  \put(128,55){\makebox(0,0){$\Vert$}}
  \put(128,25){\makebox(0,0){$\Vert$}}
 \end{picture}
 \par\smallskip
 \renewcommand{\thefigure}{4.6.2}
 \refstepcounter{figure}
 \textbf{图式~\thefigure}\label{diag:4-6-2}
\end{figure}

现在定理 4.6.25 即由引理 4.6.40 和 4.6.43 推出。
\end{proof}

至此，定理 4.1.2 已得证明。
定理 4.1.1 与 4.1.2 也有同伦含幺版本：
上述同伦等价
$\mathfrak g^{\psi,\{J_\rho\}_\rho,
\mathfrak s^{\operatorname{top}(\rho)}}$
是含幺的，而定理 4.6.25 中的“同伦”
可加强为“\mbox{u-同伦}”。
这一版本可由本节以及第 7.2、7.3 节的论证合并得到。
